% Lean: EconLean/Macro/PublicFinance/{SaezTopRate,InverseEuler}.lean, EconLean/Macro/Granularity.lean,
%       EconLean/Finance/RareDisasters.lean
\chapter{Top Taxes, the Inverse Euler Equation, Granularity and Rare Disasters}
\label{chap:public_finance_risk}

Sources: Saez (2001); Diamond--Saez (2011); Piketty--Saez (2013); Rogerson (1985); Golosov--Kocherlakota--Tsyvinski (2003);
Kocherlakota, \emph{The New Dynamic Public Finance}; Gabaix (2011); Rietz (1988); Barro (2006).

\section{The optimal top income tax rate}

\begin{definition}[Top-tax problem]
\label{def:saez}
\lean{Econ.Saez.Params, Econ.Saez.Params.zbar, Econ.Saez.Params.dSW, Econ.Saez.Params.topRate, Econ.Saez.Params.welfare}
\leanok
Incomes above $z^*>0$ are Pareto with tail $a>1$, mean $\bar z=az^*/(a-1)$; $e>0$ is the elasticity of top incomes with respect
to $1-\tau$ and $\bar g\in[0,1)$ the welfare weight. $dSW(\tau)=(1-\bar g)(\bar z-z^*)-e\bar z\tau/(1-\tau)$,
$\tau^*=(1-\bar g)/(1-\bar g+ae)$ and $W(\tau)=(1-\bar g)(\bar z-z^*)\tau+e\bar z(\tau+\log(1-\tau))$.
\end{definition}

\begin{theorem}[The Saez formula]
\label{thm:saez}
\lean{Econ.Saez.mean_pareto, Econ.Saez.Params.zbar_eq_integral, Econ.Saez.Params.zbar_sub, Econ.Saez.Params.zbar_pos, Econ.Saez.Params.den_pos, Econ.Saez.Params.topRate_mem, Econ.Saez.Params.dSW_scaled, Econ.Saez.Params.scale_pos, Econ.Saez.Params.dSW_eq_zero_iff, Econ.Saez.Params.dSW_pos_iff, Econ.Saez.Params.revenue_max, Econ.Saez.Params.topRate_anti_e, Econ.Saez.Params.topRate_anti_a, Econ.Saez.Params.welfare_hasDerivAt, Econ.Saez.Params.welfare_hasDerivAt_dSW, Econ.Saez.Params.hasDerivAt_welfare, Econ.Saez.Params.topRate_isMaxOn}
\leanok
\uses{def:saez}
If mean top income $Z$ has $Z'=-eZ/(1-\tau)$ and top taxpayers' money-metric utility has $V'=-(Z-z^*)$ (envelope theorem),
then $\tau(Z-z^*)+\bar gV$ has derivative $(1-\bar g)(Z-z^*)-eZ\tau/(1-\tau)$, which equals $dSW(\tau)$ where $Z=\bar z$.
On $[0,1)$, $dSW(\tau)>0$ iff $\tau<\tau^*$ and $dSW(\tau)=0$ iff $\tau=\tau^*$; hence $\tau^*\in(0,1)$ maximizes $W$ on
$[0,1)$. With $\bar g=0$, $\tau^*=1/(1+ae)$; $\tau^*$ falls with $e$ and with $a$.
\end{theorem}
\begin{proof}
\leanok
Product rule for the derivative. Since $\bar z-z^*=\bar z/a$, $dSW(\tau)\cdot(1-\tau)a/\bar z=(1-\bar g)-\tau(1-\bar g+ae)$ with
a positive factor. $W'=dSW$, so $W$ is monotone on $[0,\tau^*]$ and antitone on $[\tau^*,1)$ by the mean value theorem.
\end{proof}

\section{The inverse Euler equation}

\begin{definition}[Utility perturbation]
\label{def:inverse_euler}
\lean{Econ.InverseEuler.perturbCost}
\leanok
With $h=u^{-1}$, lowering today's utility by $\Delta$ and raising every state's utility tomorrow by $\Delta/\beta$ costs
$\Phi(\Delta)=h(v_0-\Delta)+R^{-1}\sum_s\pi_sh(v_s+\Delta/\beta)$.
\end{definition}

\begin{theorem}[Inverse Euler equation and savings wedge]
\label{thm:inverse_euler}
\lean{Econ.InverseEuler.utility_neutral, Econ.InverseEuler.incentive_neutral, Econ.InverseEuler.inverse_euler, Econ.InverseEuler.mean_mul_mean_inv_sub, Econ.InverseEuler.one_le_mean_mul_mean_inv, Econ.InverseEuler.one_lt_mean_mul_mean_inv, Econ.InverseEuler.savings_wedge, Econ.InverseEuler.inverse_euler_wedge}
\leanok
\uses{def:inverse_euler}
With $\sum\pi=1$ the perturbation leaves lifetime utility and all utility differences across states unchanged. If $\Delta=0$
minimizes $\Phi$ locally, then $h'(v_0)=(\beta R)^{-1}\sum\pi_sh'(v_s)$, i.e.\ $1/u'(c_0)=(\beta R)^{-1}E[1/u'(c_1)]$. Then
$u'(c_0)\le\beta RE[u'(c_1)]$, strictly if marginal utility is not constant across positive-probability states.
\end{theorem}
\begin{proof}
\leanok
$\Phi'(0)=-h'(v_0)+(\beta R)^{-1}\sum\pi_sh'(v_s)=0$. For the wedge,
$(\sum\pi x)(\sum\pi/x)-1=\tfrac12\sum_{i,j}\pi_i\pi_j(x_i-x_j)^2/(x_ix_j)\ge0$.
\end{proof}

\section{Granular fluctuations}

\begin{definition}[Aggregate growth and Zipf shares]
\label{def:granular}
\lean{Econ.Granular.agg, Econ.Granular.zipf}
\leanok
GDP growth is $\sum_is_i\varepsilon_i$. Zipf shares are $s_i=(1/i)/H_N$ with $H_N$ the harmonic number.
\end{definition}

\begin{theorem}[Granular volatility]
\label{thm:granular}
\lean{Econ.Granular.E_agg, Econ.Granular.var_agg, Econ.Granular.herfindahl_ge, Econ.Granular.herfindahl_equal, Econ.Granular.granular_bound, Econ.Granular.var_equal, Econ.Granular.harmonic_eq_sum, Econ.Granular.harmonic_pos, Econ.Granular.zipf_sum, Econ.Granular.zipf_herfindahl, Econ.Granular.var_zipf_ge}
\leanok
\uses{def:granular, def:capm}
With mean-zero uncorrelated shocks of variance $\sigma^2$, $\operatorname{Var}=\sigma^2\sum s_i^2\ge\sigma^2s_k^2$ for every $k$, and
$\sum s_i^2\ge1/N$. Equal shares give $\sigma^2/N$; Zipf shares give $\operatorname{Var}\ge\sigma^2/(1+\log N)^2$.
\end{theorem}
\begin{proof}
\leanok
Expand the square; cross terms vanish. Cauchy--Schwarz for $1/N$. For Zipf, keep the largest share $1/H_N$ and use
$H_N\le1+\log N$.
\end{proof}

\section{Rare disasters}

\begin{definition}[Disaster economy]
\label{def:disasters}
\lean{Econ.Disasters.Params, Econ.Disasters.Params.G₁, Econ.Disasters.Params.G₂, Econ.Disasters.Params.E, Econ.Disasters.Params.equityPrice, Econ.Disasters.Params.Rf, Econ.Disasters.Params.ERe}
\leanok
Growth is $g$ with probability $1-p$ and $(1-b)g$ with probability $p$; the SDF is $\beta G^{-\gamma}$; $R_f=1/E[m]$ and equity
is a claim to $G$ with expected return $E[G]/E[mG]$.
\end{definition}

\begin{theorem}[Disasters and the equity premium]
\label{thm:disasters}
\lean{Econ.Disasters.Params.G₂_pos, Econ.Disasters.Params.G₂_lt, Econ.Disasters.Params.rpow_mul_self, Econ.Disasters.Params.E_pos, Econ.Disasters.Params.premium_ratio_eq, Econ.Disasters.Params.cov_two_point, Econ.Disasters.Params.premium_pos, Econ.Disasters.Params.premium_zero, Econ.Disasters.Params.E_mu_strictMono_p, Econ.Disasters.Params.riskfree_anti_p}
\leanok
\uses{def:disasters}
$E[R_e]/R_f=E[G]E[G^{-\gamma}]/E[G^{1-\gamma}]$. With $0<p<1$ and $\gamma>0$ the premium is strictly positive; with $p=0$ it
is zero. Holding $b,g,\gamma,\beta$ fixed, a higher $p$ strictly lowers $R_f$.
\end{theorem}
\begin{proof}
\leanok
$E[G^{1-\gamma}]-E[G]E[G^{-\gamma}]=p(1-p)(G_1-G_2)(G_1^{-\gamma}-G_2^{-\gamma})<0$. A higher $p$ moves weight to the state with
larger $G^{-\gamma}$.
\end{proof}
